\documentclass[aspectratio=169,11pt]{beamer}

% ============================================================================
% THEME AND PROFESSIONAL COLORS
% ============================================================================

\usetheme{Darmstadt}
\usecolortheme{default}

% Color definitions
\definecolor{mainblue}{RGB}{50,50,120}
\definecolor{lightblue}{RGB}{240,245,255}
\definecolor{mediumblue}{RGB}{200,210,240}
\definecolor{darkblueviolet}{RGB}{40,40,100}
\definecolor{accentviolet}{RGB}{80,60,140}
\definecolor{darkgreen}{RGB}{0,100,50}
\definecolor{darkred}{RGB}{180,0,0}
\definecolor{darkblue}{RGB}{0,51,102}

% Color configuration
\setbeamercolor{structure}{fg=mainblue}
\setbeamercolor{normal text}{fg=black!90}
\setbeamercolor{frametitle}{bg=mainblue,fg=white}
\setbeamercolor{block title}{bg=mediumblue,fg=black!90}
\setbeamercolor{block body}{bg=lightblue,fg=black!85}
\setbeamercolor{block title alerted}{bg=red!20,fg=black!90}
\setbeamercolor{block body alerted}{bg=red!10,fg=black!85}

\usefonttheme{structurebold}
\setbeamercovered{invisible}
\setbeamertemplate{navigation symbols}{}
\setbeamertemplate{footline}[frame number]

% Custom frame template to prevent overflows
\setbeamertemplate{frametitle}{%
	\nointerlineskip%
	\tikz[overlay,remember picture] {%
		\fill[mainblue] (current page.north west) rectangle ([yshift=-4.125ex]current page.north east);
	}%
	\begin{beamercolorbox}[wd=\paperwidth,ht=2.75ex,dp=1.375ex,leftskip=0.3cm,rightskip=0.3cm,sep=0pt]{frametitle}%
		\usebeamerfont{frametitle}%
		\color{white}%
		\vbox{}\vskip-2ex%
		\strut\insertframetitle\strut%
		\vskip-0.5ex%
	\end{beamercolorbox}%
	\vskip0.3em%
}

% Ensure frametitle has correct background color (after template definition)
\setbeamercolor{frametitle}{bg=mainblue,fg=white}

% ============================================================================
% PACKAGES
% ============================================================================

\usepackage[utf8]{inputenc}
\usepackage[T1]{fontenc}
\usepackage{lmodern}
\usepackage[english]{babel}

\usepackage{algorithmic}
\usepackage[ruled,vlined]{algorithm2e}

% Mathematics
\usepackage{amsmath,amssymb,amsfonts,mathtools}
\usepackage{xcolor}

% Graphics
\usepackage{graphicx}
\usepackage{tikz}
\usetikzlibrary{shapes,arrows,positioning,calc}
\usepackage{adjustbox}

% Algorithms
\usepackage{algorithm}
\usepackage{algpseudocode}
\usepackage{listings}

% Symbols
\usepackage{pifont}

% Links
\usepackage{hyperref}

% ============================================================================
% LISTINGS CONFIGURATION
% ============================================================================

\lstset{
	basicstyle=\ttfamily\small,
	breaklines=true,
	keywordstyle=\color{darkblue}\bfseries,
	commentstyle=\color{gray},
	stringstyle=\color{darkred},
	showstringspaces=false,
	language=Python,
	numbers=left,
	numberstyle=\tiny\color{gray},
	frame=single,
	framexleftmargin=12pt,
	xleftmargin=12pt,
	aboveskip=5pt,
	belowskip=5pt
}

% ============================================================================
% TITLE
% ============================================================================

\title[DOA-DQN]{Application of DOA Algorithm to DQN Optimization}
\subtitle{Hyperparameter Optimization for Deep Q-Networks using Dream Optimization Algorithm}
\author{Radhi Badache}
\institute{Research Project -- Statistics for Data Science}
\date{\today}

\begin{document}
	
	% ============================================================================
	% TITLE PAGE
	% ============================================================================
	
	\begin{frame}[plain]
		\setbeamercolor{background canvas}{bg=mainblue}
		\titlepage
	\end{frame}
	
	% ============================================================================
	% TABLE OF CONTENTS
	% ============================================================================
	
	\begin{frame}{Table of Contents}
		\tableofcontents
	\end{frame}
	
	
	
	% ============================================================================
	% SECTION 2: OPTIMIZER FAMILIES
	% ============================================================================
	
	\section{Two Families of Optimizers}
	
	\begin{frame}[plain]
		\setbeamercolor{background canvas}{bg=mainblue}
		\centering
		\vspace{2.5cm}
		{\Huge\textcolor{bg=mainblue}{\textbf{Optimizer Families \\}}}
		\vspace{0.8cm}
		{\Large\textcolor{bg=mainblue}{HPO vs Parameter Optimizers}}
		\normalsize
	\end{frame}
	
	\begin{frame}{Hyperparameter vs Parameters}
		\begin{columns}[T]
			\column{0.48\textwidth}
			\begin{block}{Parameter Optimization}
				\footnotesize
				\begin{itemize}
					\item \textbf{Goal:} Find $\theta^*$ minimizing loss
					\item \textbf{Methods:} Adam, SGD, RMSprop
					\item \textbf{Context:} Training inner loop
				\end{itemize}
			\end{block}
			
			\column{0.48\textwidth}
			\begin{block}{Hyperparameter Optimization}
				\footnotesize
				\begin{itemize}
					\item \textbf{Goal:} Find $\lambda^*$ maximizing performance
					\item \textbf{Methods:} Bayesian, TPE, Population-based
					\item \textbf{Context:} Project-level tuning
				\end{itemize}
			\end{block}
		\end{columns}
	\end{frame}
	
	\begin{frame}{Hyperparameters to Optimize}
		\begin{columns}[T]
			\column{0.48\textwidth}
			\begin{block}{DQN Training Parameters}
				\footnotesize
				\begin{itemize}
					\item \textbf{learning\_rate} (float)
					\item \textbf{gamma} (float): discount factor
					\item \textbf{epsilon\_start} (float)
					\item \textbf{epsilon\_end} (float)
					\item \textbf{epsilon\_decay\_steps} (int)
					\item \textbf{batch\_size} (int)
					\item \textbf{replay\_buffer\_size} (int)
				\end{itemize}
			\end{block}
			
			\column{0.48\textwidth}
			\begin{block}{Network \& Optimizer}
				\footnotesize
				\begin{itemize}
					\item \textbf{target\_update\_freq} (int)
					\item \textbf{hidden\_size} (int)
					\item \textbf{n\_layers} (int)
					\item \textbf{adam\_beta1} (float)
					\item \textbf{adam\_beta2} (float)
					\item \textbf{adam\_epsilon} (float)
				\end{itemize}
			\end{block}
		\end{columns}
		
		\vspace{0.3cm}
		
		\begin{block}{Total: 13 Hyperparameters}
			\footnotesize
			DOA algorithm optimizes all these hyperparameters simultaneously to maximize DQN performance on Atari games.
		\end{block}
	\end{frame}

	\begin{frame}{DOA Algorithm Flow}

	\begin{center}
		\begin{tikzpicture}[
			node distance = 1.2cm,
			auto,
			thick,
			block/.style={rectangle, draw=mainblue, fill=lightblue, text width=2.2cm, text centered, rounded corners, minimum height=0.7cm, font=\footnotesize},
			decision/.style={diamond, draw=mainblue, fill=mediumblue, text width=1.5cm, text centered, minimum height=0.8cm, font=\footnotesize},
			arrow/.style={->, >=stealth, mainblue, thick}
			]
			% Start
			\node [block] (init) {Initialize};
			
			% Evaluate
			\node [block, below of=init] (eval) {Evaluate};
			
			% Decision: exploration or exploitation
			\node [decision, below of=eval, yshift=-0.2cm] (decision) {$t < T_d$?};
			
			% Exploration branch
			\node [block, below left of=decision, xshift=-1.2cm, yshift=-0.4cm, text width=1.8cm] (explore) {Exploration\\5 Groups};
			
			% Exploitation branch
			\node [block, below right of=decision, xshift=1.2cm, yshift=-0.4cm, text width=1.8cm] (exploit) {Exploitation\\Global Best};
			
			% Merge
			\node [block, below of=decision, yshift=-2cm, text width=2.2cm] (update) {Update\\Elitism};
			
			% Loop back
			\node [block, left of=eval, xshift=-3cm] (loop) {Iterate};
			
			% End
			\node [block, right of=update, xshift=3cm] (end) {Best Solution};
			
			% Arrows
			\draw [arrow] (init) -- (eval);
			\draw [arrow] (eval) -- (decision);
			\draw [arrow] (decision) -- node[left, font=\tiny] {Yes} (explore);
			\draw [arrow] (decision) -- node[right, font=\tiny] {No} (exploit);
			\draw [arrow] (explore) |- (update);
			\draw [arrow] (exploit) |- (update);
			\draw [arrow] (update) -| node[below, font=\tiny, yshift=-0.2cm] {$t < T_{\max}$} (loop);
			\draw [arrow] (loop) -- (eval);
			\draw [arrow] (update) -| node[above, font=\tiny, yshift=0.2cm] {$t = T_{\max}$} (end);
		\end{tikzpicture}
	\end{center}
	
	\vspace{0.1cm}
	
	\begin{block}{Key Phases}
		\footnotesize
		\textbf{Exploration} ($t < 0.9 \times T_{\max}$): 5 groups with memory, forgetting, and dream-sharing. \textbf{Exploitation} ($t \geq 0.9 \times T_{\max}$): Refinement around global best.
	\end{block}
\end{frame}


% ============================================================================
% REFERENCES
% ============================================================================

\section{References}

\begin{frame}[fragile,allowframebreaks]{\small References}
	\begin{thebibliography}{99}
		
		\bibitem{Lang2025} Lang, Y., Gao, Y. (2025). Dream Optimization Algorithm (DOA): A novel metaheuristic optimization algorithm. \textit{Computer Methods in Applied Mechanics and Engineering}, 436, 117718.
		
		\bibitem{Mnih2013} Mnih, V., Kavukcuoglu, K., Silver, D. (2013). Playing Atari with Deep Reinforcement Learning. \textit{arXiv preprint arXiv:1312.5602}.
		
		\bibitem{Mnih2015} Mnih, V., Kavukcuoglu, K., Silver, D., et al. (2015). Human-level control through deep reinforcement learning. \textit{Nature}, 529(7587), 529-533.
		
		\bibitem{Watkins1992} Watkins, C. J. C. H., Dayan, P. (1992). Q-learning. \textit{Machine Learning}, 8(3), 279-292.
		
		\bibitem{Kennedy1995} Kennedy, J., Eberhart, R. (1995). Particle swarm optimization. \textit{Proceedings of ICNN'95}, 4, 1942-1948.
		
		\bibitem{Mirjalili2014} Mirjalili, S., Mirjalili, S. M., Lewis, A. (2014). Grey Wolf Optimizer. \textit{Advances in Engineering Software}, 69, 46-61.
		
		\bibitem{Mirjalili2016} Mirjalili, S., Lewis, A. (2016). The Whale Optimization Algorithm. \textit{Advances in Engineering Software}, 95, 51-67.
		
		\bibitem{Bergstra2011} Bergstra, J., Bengio, Y. (2011). Hyperparameter Optimization with Approximate Gradient. \textit{ICML}, 24, 265-273.
		
		\bibitem{Sutton1998} Sutton, R. S., Barto, A. G. (1998). \textit{Reinforcement Learning: An Introduction}. MIT Press.
		
	\end{thebibliography}
\end{frame}

\end{document}